\subsection{Передаточная функция обратной связи}

В зависимости от структурной схемы цепи обратной связи формируется уравнение сигнала управления. Исходная схема представлена на рисунке~\ref{fig:feedback_scheme}.

\begin{figure}[H]
    \centering
    \centering
\begin{tikzpicture}[
    scale=0.79,
    transform shape,
    node distance=1.5cm,
    >=Stealth,
    block/.style={
        draw,
        rectangle,
        minimum width=2.8cm,
        minimum height=1.1cm,
        align=center
    },
    smallblock/.style={
        draw,
        rectangle,
        minimum width=1.25cm,
        minimum height=0.9cm,
        align=center
    },
    sum/.style={
        draw,
        circle,
        minimum size=0.72cm,
        inner sep=0pt,
        path picture={
            \draw[thick]
                (path picture bounding box.south west) --
                (path picture bounding box.north east);
            \draw[thick]
                (path picture bounding box.north west) --
                (path picture bounding box.south east);
        }
    },
    line/.style={draw, -Stealth, thick}
]

    % Основная цепь
    \node[coordinate] (input) {};
    \node[block, right=1.4cm of input] (motor) {Двигатель};
    \node[block, right=1.8cm of motor] (gear) {Редуктор};
    \node[block, right=1.8cm of gear] (load) {Выходное\\звено};

    \draw[line] (input) -- node[above] {$u$} (motor);
    \draw[line] (motor) -- node[above] {$\varphi_{\text{р}}$} (gear);
    \draw[line] (gear) -- node[above] {$\varphi_{\text{м}}$} (load);
    \draw[line] (load) -- ++(2,0) node[right] {$\varphi_{\text{м}}$};

    % Измерительные звенья
    \node[smallblock, below=1.5cm of motor] (kr) {$\dfrac{k}{i}$};
    \node[smallblock, below=2.7cm of load] (km) {$k$};

    \draw[line] ($(motor.east)!0.35!(gear.west)$) |- ++(0,-1.5) -| (kr.north);
    \draw[line] ($(gear.east)!0.35!(load.west)$) |- ++(0,-1.5) -| (km.north);

    % Задающее воздействие
    \node[block, below=1.5cm of gear] (ref) {$k\varphi^*$};

    % Сумматоры рассогласования
    \node[sum, below=2.3cm of kr] (sumr) {};
    \node[sum, below=1.5cm of km] (summ) {};

    % Закрашенные сектора сумматоров рассогласования
    \fill[pattern=north east lines]
        (sumr.center) -- ++(-45:0.36cm)
        arc[start angle=-45,end angle=45,radius=0.36cm] -- cycle;

    \fill[pattern=north east lines]
        (summ.center) -- ++(135:0.36cm)
        arc[start angle=135,end angle=225,radius=0.36cm] -- cycle;

    \draw[line] (kr.south) -- node[right] {$q_{\text{р}}$} (sumr.north);
    \draw[line] (km.south) -- node[right] {$q_{\text{м}}$} (summ.north);

    \draw[line] (ref.south) |- (sumr.east);
    \draw[line] (ref.south) |- (summ.west);

    % Звенья регулятора
    \node[block, left=1.5cm of sumr] (kp)
        {$W_{\text{р}}(p)=k_{\text{п}}$};

    \node[block, below=0.6cm of kp] (ki)
        {$W_{\text{м}}(p)=\dfrac{k_{\text{и}}}{p}$};

    \draw[line] (sumr.west) -- node[above] {$k\psi_{\text{р}}$} (kp.east);
    \draw[line] (summ.south) |- node[below] {$k\psi_{\text{м}}$} (ki.east);

    % Сумматор управляющего сигнала
    \node[sum, left=1.5cm of kp] (sumu) {};

    % Закрашенный сектор сумматора управляющего сигнала
    \fill[pattern=north east lines]
        (sumu.center) -- ++(45:0.36cm)
        arc[start angle=45,end angle=135,radius=0.36cm] -- cycle;

    \fill[pattern=north east lines]
        (sumu.center) -- ++(-135:0.36cm)
        arc[start angle=-135,end angle=-45,radius=0.36cm] -- cycle;

    \draw[line] (kp.west) -| ++(-.5, 1) -| (sumu.north);
    \draw[line] (ki.west) -| (sumu.south);

    % Возврат сигнала управления на двигатель
    \draw[draw, thick] (sumu.west) -- ++(-0.9,0) |- (input);

\end{tikzpicture}

    \caption{Схема обратной связи}
    \label{fig:feedback_scheme}
\end{figure}

По схеме видно, что на регулятор поступают сигналы рассогласования
\begin{equation}
    k\psi_{\text{р}} = q_{\text{р}} - q^*
    =
    k\left(
        \varphi_{\text{рпр}} - \varphi^*
    \right),
\end{equation}

\begin{equation}
    k\psi_{\text{м}} = q_{\text{м}} - q^*
    =
    k\left(
        \varphi_{\text{м}} - \varphi^*
    \right).
\end{equation}

В соответствии с вариантом используются пропорциональное звено по сигналу $\psi_{\text{р}}$ и интегральное звено по сигналу $\psi_{\text{м}}$:
\begin{equation}
    W_{\text{р}}(p) = k_{\text{п}},
    \qquad
    W_{\text{м}}(p) = \frac{k_{\text{и}}}{p}.
\end{equation}

Так как обратная связь является отрицательной, сигнал управления имеет вид
\begin{equation}
    u
    =
    -k_{\text{п}} k
    \left(
        \varphi_{\text{рпр}} - \varphi^*
    \right)
    -
    \frac{k_{\text{и}}}{p} k
    \left(
        \varphi_{\text{м}} - \varphi^*
    \right).
\end{equation}

Или, в операторной форме:
\begin{equation}
    U
    =
    -k_{\text{п}} k
    \left(
        \Phi_{\text{рпр}} - \Phi^*
    \right)
    -
    \frac{k_{\text{и}}}{p} k
    \left(
        \Phi_{\text{м}} - \Phi^*
    \right).
\end{equation}

Умножим обе части уравнения на $p$:
\begin{equation}
    pU
    =
    -k_{\text{п}}kp
    \left(
        \Phi_{\text{рпр}} - \Phi^*
    \right)
    -
    k_{\text{и}}k
    \left(
        \Phi_{\text{м}} - \Phi^*
    \right).
\end{equation}

Раскроем скобки:
\begin{equation}
    pU
    =
    -k_{\text{п}}kp\Phi_{\text{рпр}}
    -
    k_{\text{и}}k\Phi_{\text{м}}
    +
    k_{\text{п}}kp\Phi^*
    +
    k_{\text{и}}k\Phi^*.
\end{equation}

Так как в дальнейшем рассматривается постоянное задающее воздействие $\varphi^* = \operatorname{const}$, то $\dot{\varphi}^* = 0$. Тогда после перехода от оператора $p$ к производной по времени получим:
\begin{equation}
    \dot{u}
    =
    -k_{\text{п}}k\dot{\varphi}_{\text{рпр}}
    -
    k_{\text{и}}k\varphi_{\text{м}}
    +
    k_{\text{и}}k\varphi^*.
\end{equation}

Итоговое уравнение системы управления запишем в виде
\begin{empheq}[box=\fbox]{equation}
    \dot{u}
    +
    k_{\text{п}}k\dot{\varphi}_{\text{рпр}}
    +
    k_{\text{и}}k\varphi_{\text{м}}
    =
    k_{\text{и}}k\varphi^*.
\end{empheq}
